% BHU PYQ -- Manually transcribed & error-corrected
% Paper: BCF-123 : Fundamental of Accounting
% Exam: B.Com. (Hons.) FMM Semester II Examination, 2021-22

\documentclass[11pt,a4paper]{article}
\usepackage[a4paper,top=2.5cm,bottom=2.5cm,left=2.8cm,right=2.8cm,headheight=14pt]{geometry}
\usepackage{amsmath,amssymb}
\usepackage[T1]{fontenc}
\usepackage[utf8]{inputenc}
\usepackage{lmodern,microtype,fancyhdr,enumitem,hyperref,lastpage,parskip}
\usepackage{booktabs}

\hypersetup{
    colorlinks=true,
    urlcolor=black,
    linkcolor=black,
    pdftitle={BCF-123: Fundamental of Accounting -- BHU B.Com. (Hons.) FMM Sem II 2021-22}
}

\pagestyle{fancy}
\fancyhf{}
\renewcommand{\headrulewidth}{0.4pt}
\renewcommand{\footrulewidth}{0.4pt}
\fancyhead[L]{\small   \textit{Banaras Hindu University}}
\fancyhead[R]{\small   \textit{BCF-123: Fundamental of Accounting}}
\fancyfoot[C]{\small Page  \thepage\ of \pageref{LastPage}}


\newlist{parts}{enumerate}{2}
\setlist[parts,1]{label=(\alph*),leftmargin=*,itemsep=5pt,topsep=3pt}
\setlist[parts,2]{label=(\roman*),leftmargin=*,itemsep=3pt,topsep=2pt}

\newcommand{\pts}[1]{\hfill   \textit{[#1]}}


\begin{document}


\begin{center}
  {\large\bfseries BANARAS HINDU UNIVERSITY}\\[2pt]
  {\normalsize B.Com. (Hons.) FMM --- Semester II Examination, 2021-22}\\[6pt]
  \rule{0.8\linewidth}{0.5pt}\\[6pt]
  {\large\bfseries Paper No. BCF-123: Fundamental of Accounting}\\[4pt]
  {\normalsize Time: 3 Hours\hfill Maximum Marks: 70}
\end{center}

\rule{\linewidth}{0.4pt}

\smallskip



\noindent   \textit{Instructions: Answer all questions. All questions carry equal marks (14 marks each).}

\bigskip




\noindent   \textbf{Question 1.}
Why is a distinction made between capital and revenue expenditure? Explain by taking an imaginary illustration.\pts{14}

\centerline{   \textbf{OR}}

On April 1, 2021, A started business with \pounds 1,000,000 (read as Rupees) and other transactions for the month are as follows:

\begin{itemize}[itemsep=2pt,topsep=2pt]
  \item[April 2:] Purchased furniture for cash \pounds 7,000.
  \item[April 8:] Purchased goods for cash \pounds 2,000 and on credit \pounds 1,000 from K Retail Store.
  \item[April 14:] Sold goods to Khan Brothers \pounds 12,000 and cash sales \pounds 5,000.
  \item[April 18:] Owner withdrew goods worth \pounds 2,000 for personal use.
  \item[April 22:] Paid K Retail Store \pounds 500.
  \item[April 26:] Received \pounds 10,000 from Khan Brothers.
  \item[April 27:] Gave away charity of cash \pounds 50 and merchandising (goods) worth \pounds 30.
  \item[April 28:] Bad debts during the period were \pounds 100.
  \item[April 29:] Payment made to creditors to the value of \pounds 20,000 at 10 percent discount.
  \item[April 30:] Wages paid \pounds 1,200 (including \pounds 200 relating to a future year).
\end{itemize}
Pass the Journal Entries in the books of A.\pts{14}

\medskip\rule{\linewidth}{0.2pt}


\newpage




\noindent   \textbf{Question 2.}
The Trial Balance of a business as at 31st March, 2021 is given below:


\begin{center}

\begin{tabular}{lr|lr}
 oprule
   \textbf{Debit Balances} &   \textbf{Amount (\pounds)} &   \textbf{Credit Balances} &   \textbf{Amount (\pounds)} \\
\midrule
Stock (1 April, 2020) & 25,000 & Sales & 2,27,800 \\
Furniture & 8,000 & Commission & 500 \\
Plant and Machinery & 1,50,000 & Returns Outward & 1,000 \\
Debtors & 30,000 & Creditors & 40,000 \\
Wages & 12,000 & Capital & 1,50,000 \\
Salaries & 20,000 & & \\
Bad Debts & 1,000 & & \\
Purchases & 1,20,000 & & \\
Electricity Charges & 1,200 & & \\
Telephone Charges & 2,400 & & \\
General Expenses & 3,000 & & \\
Postage Expenses & 1,800 & & \\
Returns Inward & 900 & & \\
Insurance Premium & 1,500 & & \\
Cash in Hand & 2,500 & & \\
Cash at Bank & 40,000 & & \\
\midrule
   \textbf{Total} &   \textbf{4,19,300} &   \textbf{Total} &   \textbf{4,19,300} \\
ottomrule
\end{tabular}
\end{center}




\noindent   \textbf{Adjustments:}

\begin{parts}
  \item Closing stock was valued at \pounds 7,000.
  \item Outstanding liabilities for wages were \pounds 600 and salaries \pounds 1,400.
  \item Depreciation is to be provided @ 5\% p.a. on all fixed assets.
  \item Included in plant and machinery is a machine purchased for \pounds 10,000 on 1st October, 2020.
  \item Insurance premium paid in advance was \pounds 200.
\end{parts}
Prepare Trading and Profit \& Loss Account for the year ended 31st March, 2021 and a Balance Sheet as at that date.\pts{14}

\centerline{   \textbf{OR}}

What is a Bank Reconciliation Statement? Explain the reasons where the balance shown by the Bank Pass Book does not agree with the balance as shown by the bank column of the Cash Book.\pts{14}

\medskip\rule{\linewidth}{0.2pt}


\newpage




\noindent   \textbf{Question 3.}

\begin{parts}
  \item Radha and Rukmani are partners in a firm sharing profits in the ratio of 3:2. They admitted Gopi as a new partner. Radha surrendered 1/3 of her share in favour of Gopi and Rukmani surrendered 1/4 of her share in favour of Gopi. Calculate the new profit sharing ratio.\pts{5}
  \item A and B are partners in a firm sharing profits and losses in the ratio of 3:2. They decided to admit C into partnership with a 1/4 share in profits. C will bring in \pounds 50,000 as his capital and the requisite amount of goodwill premium in cash. The goodwill of the firm is valued at \pounds 1,20,000. The new profit sharing ratio is agreed to be 5:4:1. A and B withdrew their share of goodwill. Give the necessary journal entries.\pts{9}
\end{parts}

\centerline{   \textbf{OR}}

Elaborate the term `Goodwill'. Explain the methods of valuation of goodwill with illustration.\pts{14}

\medskip\rule{\linewidth}{0.2pt}




\noindent   \textbf{Question 4.}
Given below is the Balance Sheet of A and B, who are carrying on a partnership business as on 31.12.2021. A and B share profits and losses in the ratio of 2:1:


\begin{center}

\begin{tabular}{lr|lr}
 oprule
   \textbf{Liabilities} &   \textbf{Amount (\pounds)} &   \textbf{Assets} &   \textbf{Amount (\pounds)} \\
\midrule
Bills Payable & 10,000 & Cash in Hand & 10,000 \\
Creditors & 58,000 & Cash at Bank & 40,000 \\
Outstanding Expenses & 2,000 & Sundry Debtors & 60,000 \\
Capital Accounts: & & Stock & 40,000 \\
~~A & 1,80,000 & Plant & 1,00,000 \\
~~B & 1,50,000 & Building & 1,50,000 \\
\midrule
   \textbf{Total} &   \textbf{4,00,000} &   \textbf{Total} &   \textbf{4,00,000} \\
ottomrule
\end{tabular}
\end{center}




\noindent C is admitted as a partner on the date of the Balance Sheet on the following terms:

\begin{parts}
  \item C will bring in \pounds 1,00,000 as his capital and \pounds 60,000 as his share of goodwill for a 1/4 share in the profits.
  \item Plant is to be appreciated to \pounds 1,20,000 and the value of buildings is to be appreciated by 10\%.
  \item Stock is found over-valued by \pounds 4,000.
  \item A provision for bad and doubtful debts is to be created at 5\% of debtors.
  \item Creditors were unrecorded to the extent of \pounds 1,000.
\end{parts}
Prepare the Revaluation Account, Partners' Capital Accounts, and show the Balance Sheet after the admission of C.\pts{14}

\centerline{   \textbf{OR}}

Do you advise that assets and liabilities must be revalued at the time of admission of a partner? If so, why? Also describe how this is treated in the books of accounts.\pts{14}

\medskip\rule{\linewidth}{0.2pt}


\newpage




\noindent   \textbf{Question 5.}
How is the Joint Life Insurance Policy recorded in the books of a partnership firm at the time of retirement of a partner? Discuss the alternative methods.\pts{14}

\centerline{   \textbf{OR}}

Dhaval, Kamal, and Naval are partners sharing profits and losses in the ratio of 2:2:1. Naval retires on 31.03.2021. The Balance Sheet of the firm as on 31.03.2021 was as under:


\begin{center}

\begin{tabular}{lr|lr}
 oprule
   \textbf{Liabilities} &   \textbf{Amount (\pounds)} &   \textbf{Assets} &   \textbf{Amount (\pounds)} \\
\midrule
Capital Accounts: & & Goodwill & 10,000 \\
~~Dhaval & 30,000 & Machinery & 20,000 \\
~~Kamal & 20,000 & Investment & 10,000 \\
~~Naval & 10,000 & Debtors & 30,000 \\
General Reserve & 5,000 & Stock & 10,000 \\
Investment Fluctuation Fund & 2,500 & Cash at Bank & 5,000 \\
Bad Debts Reserve & 2,000 & & \\
Creditors & 15,500 & & \\
\midrule
   \textbf{Total} &   \textbf{85,000} &   \textbf{Total} &   \textbf{85,000} \\
ottomrule
\end{tabular}
\end{center}




\noindent The following adjustments are agreed upon at the time of retirement:

\begin{parts}
  \item The value of machinery will be \pounds 25,000 and the value of stock is \pounds 5,000.
  \item The value of investments is \pounds 8,000, which is taken by Naval at this price.
  \item An amount of \pounds 5,000 included in creditors is no longer payable.
  \item The provision for workmen compensation is to be created at \pounds 2,000.
  \item The provision for doubtful debts is to be kept at 10\% on debtors.
  \item The goodwill of the firm is valued at \pounds 40,000.
\end{parts}
Prepare the Revaluation Account, Partners' Capital Accounts, and the Balance Sheet of the firm after Naval's retirement.\pts{14}

\vfill
\rule{\linewidth}{0.4pt}

\begin{center}\small
\href{https://bhu-exam.vercel.app/}{Click here to visit our website}\\[4pt]
\href{https://me-aryan.vercel.app/}{Click here to visit our contributor}\\[4pt]
\href{https://quashan-me.vercel.app/}{Click here to visit our contributor}
\end{center}

\end{document}
